\begin{tomtat}
	\begin{minipage}[t]{0.3\textwidth}
		\vspace{0pt}
		\begin{tomtat}
			\begin{center}
				\textbf{ĐỀ 54}
			\end{center}
			
		\end{tomtat}
	\end{minipage}
	\hfill
	\begin{minipage}[t]{0.5\textwidth}
		\vspace{0pt}
		\begin{tomtat}
			\begin{center}
				\includegraphics[scale=0.08]{images/54}
			\end{center}
		\end{tomtat}
	\end{minipage}
\end{tomtat}
\textbf{I. PHẦN TRẮC NGHIỆM}
\begin{ex}
	Điều kiện xác định biểu thức $\sqrt{3-7x}$ xác định là
	\choice{$x \ge \dfrac{3}{7}$}{$x \le \dfrac{3}{7}$}{$x \le \dfrac{7}{3}$}{$x \ge \dfrac{7}{3}$}
	\loigiai{
		\textbf{Đáp án đúng là B.} \\
		Biểu thức $\sqrt{3-7x}$ xác định khi và chỉ khi $3-7x \ge 0$.
		\begin{align*}
			3-7x &\ge 0 \\
			-7x &\ge -3 \\
			x &\le \dfrac{-3}{-7} \\
			x &\le \dfrac{3}{7}
		\end{align*}
		Vậy, điều kiện xác định của biểu thức là $x \le \dfrac{3}{7}$.
	}
\end{ex}

\begin{ex}
	Biểu thức $\sqrt[3]{(\sqrt{2}-3)^3}$ có giá trị là:
	\choice{$\sqrt{2}-3$}{$3-\sqrt{2}$}{$|3-\sqrt{2}|$}{$|\sqrt{2}-3|$}
	\loigiai{
		\textbf{Đáp án đúng là A.} \\
		Ta có tính chất $\sqrt[3]{a^3} = a$ với mọi số thực $a$.
		Do đó, $\sqrt[3]{(\sqrt{2}-3)^3} = \sqrt{2}-3$.
	}
\end{ex}

\begin{ex}
	Nghiệm của bất phương trình $3x-9 \ge 0$ là:
	\choice{$x < 3$}{$x \le 3$}{$x \ge 3$}{$x \ge -3$}
	\loigiai{
		\textbf{Đáp án đúng là C.} \\
		Ta có:
		\begin{align*}
			3x-9 &\ge 0 \\
			3x &\ge 9 \\
			x &\ge \dfrac{9}{3} \\
			x &\ge 3
		\end{align*}
		Vậy, nghiệm của bất phương trình là $x \ge 3$.
	}
\end{ex}

\begin{ex}
	Gieo một con xúc xắc 50 lần cho kết quả như sau:
	\begin{center}
			\begin{tabular}{|c|c|c|c|c|c|c|}
			\hline
			Số chấm xuất hiện & 1 & 2 & 3 & 4 & 5 & 6 \\
			\hline
			Tần số & 8 & 7 & ? & 8 & 6 & 11 \\
			\hline
		\end{tabular}
	\end{center}

	Tần số xuất hiện của mặt 3 chấm là
	\choice{9}{10}{11}{12}
	\loigiai{
		\textbf{Đáp án đúng là B.} \\
		Tổng số lần gieo xúc xắc là 50 lần.
		Tổng tần số của các mặt đã biết là: $8 + 7 + 8 + 6 + 11 = 40$.
		Gọi tần số xuất hiện của mặt 3 chấm là $x$.
		Ta có: $8 + 7 + x + 8 + 6 + 11 = 50$
		\begin{align*}
			40 + x &= 50 \\
			x &= 50 - 40 \\
			x &= 10
		\end{align*}
		Vậy, tần số xuất hiện của mặt 3 chấm là 10.
	}
\end{ex}
\begin{ex}
Điểm M(-1;2) thuộc đồ thị hàm số $y=ax^2$ khi a bằng
\choice{$a=2$}{$a=4$}{$a=-4$}{$a=-2$}
\loigiai{
	\textbf{Đáp án đúng là A.} \\
	Điểm M(-1;2) thuộc đồ thị hàm số $y=ax^2$ nên tọa độ điểm M thỏa mãn phương trình hàm số:
	$2 = a(-1)^2$
	\begin{align*}
		2 &= a \times 1 \\
		a &= 2
	\end{align*}
	Vậy, $a=2$.
}
\end{ex}

\begin{ex}
Gieo một con xúc xắc cân đối và đồng chất. Gọi biến cố A: “số chấm xuất hiện trên con xúc xắc là số nguyên tố”. Số kết quả thuận lợi cho biến cố A là:
\choice{6}{2}{3}{4}
\loigiai{
	\textbf{Đáp án đúng là C.} \\
	Các số chấm có thể xuất hiện khi gieo một con xúc xắc là: $\{1, 2, 3, 4, 5, 6\}$.
	Các số nguyên tố trong tập hợp này là: $\{2, 3, 5\}$.
	Vậy, có 3 kết quả thuận lợi cho biến cố A.
}
\end{ex}

\begin{ex}
Công thức tính thể tích V của hình trụ có bán kính đáy r, chiều cao h là
\choice{$V=\dfrac{1}{3}\pi r^2h$}{$V=\pi r^2h$}{$V=\dfrac{1}{3}\pi rh^2$}{$V=\pi rh^2$}
\loigiai{
	\textbf{Đáp án đúng là B.} \\
	Công thức tính thể tích V của hình trụ có bán kính đáy r và chiều cao h là: $V = \pi r^2h$.
}
\end{ex}

\begin{ex}
Cho (O;R). sđ $\stackrel\frown{MaN} = 120^\circ$; diện tích hình quạt tròn OMaN bằng:
Hãy chọn kết quả đúng.
\choice{$\dfrac{2\pi R}{3}$}{$\dfrac{\pi R^2}{6}$}{$\dfrac{\pi R^2}{4}$}{$\dfrac{\pi R^2}{3}$}
\loigiai{
	\textbf{Đáp án đúng là D.} \\
	Diện tích hình quạt tròn được tính theo công thức: $S_q = \dfrac{\pi R^2 n}{360}$, trong đó n là số đo cung của hình quạt.
	Với $n = 120^\circ$, ta có:
	$S_q = \dfrac{\pi R^2 \times 120}{360} = \dfrac{\pi R^2}{3}$.
}
\end{ex}
\newpage
\textbf{II. PHẦN TỰ LUẬN}
\begin{ex}
	Rút gọn biểu thức: $P = \left(\dfrac{2}{\sqrt{x}+1} - \dfrac{\sqrt{x}-2}{x+\sqrt{x}}\right) : \dfrac{2}{\sqrt{x}+1}$ với $x>0$.
	\loigiai{
		Với $x > 0$, ta có:
		\begin{align*} P &= \left(\dfrac{2}{\sqrt{x}+1} - \dfrac{\sqrt{x}-2}{\sqrt{x}(\sqrt{x}+1)}\right) : \dfrac{2}{\sqrt{x}+1} \\ &= \left(\dfrac{2\sqrt{x}}{{\sqrt{x}(\sqrt{x}+1)}} - \dfrac{\sqrt{x}-2}{\sqrt{x}(\sqrt{x}+1)}\right) \cdot \dfrac{\sqrt{x}+1}{2} \\ &= \dfrac{2\sqrt{x} - (\sqrt{x}-2)}{\sqrt{x}(\sqrt{x}+1)} \cdot \dfrac{\sqrt{x}+1}{2} \\ &= \dfrac{2\sqrt{x} - \sqrt{x} + 2}{\sqrt{x}(\sqrt{x}+1)} \cdot \dfrac{\sqrt{x}+1}{2} \\ &= \dfrac{\sqrt{x}+2}{\sqrt{x}(\sqrt{x}+1)} \cdot \dfrac{\sqrt{x}+1}{2} \\ &= \dfrac{\sqrt{x}+2}{2\sqrt{x}}\end{align*}
		Vậy, $P = \dfrac{\sqrt{x}+2}{2\sqrt{x}}$ với $x>0$.
	}
\end{ex}

\begin{ex}
	Biết rằng phương trình bậc hai $x^2+4x-3=0$ có hai nghiệm phân biệt $x_1, x_2$. Không giải phương trình, tính giá trị của biểu thức $R = x_1^2x_2 + x_1x_2^2 - x_1 - x_2$.
	\loigiai{
		Phương trình $x^2+4x-3=0$ có $a=1, b=4, c=-3$.
		Do $ac = 1 \cdot (-3) = -3 < 0$ nên phương trình luôn có hai nghiệm phân biệt $x_1, x_2$.
		Theo hệ thức Vi-ét, ta có:
		$\begin{cases} x_1 + x_2 = -\dfrac{b}{a} = -\dfrac{4}{1} = -4 \\ x_1x_2 = \dfrac{c}{a} = \dfrac{-3}{1} = -3 \end{cases}$
		Ta có:
		\begin{align*} R &= x_1^2x_2 + x_1x_2^2 - x_1 - x_2 \\ &= x_1x_2(x_1+x_2) - (x_1+x_2) \\ &= (x_1x_2 - 1)(x_1+x_2) \end{align*}
		Thay giá trị $x_1+x_2 = -4$ và $x_1x_2 = -3$ vào biểu thức R:
		$R = (-3 - 1)(-4) = (-4)(-4) = 16$.
		Vậy, $R=16$.
	}
\end{ex}

\begin{ex}
	Giải hệ phương trình: $\begin{cases} x+2y=5 \\ 3x-y=1 \end{cases}$
	\loigiai{
		Ta có hệ phương trình:
		$\begin{cases} x+2y=5 \quad (1) \\ 3x-y=1 \quad (2) \end{cases}$
		Từ phương trình (2), ta suy ra $y = 3x-1$.
		Thế $y = 3x-1$ vào phương trình (1):
		\begin{align*} x + 2(3x-1) &= 5 \\ x + 6x - 2 &= 5 \\ 7x &= 5+2 \\ 7x &= 7 \\ x &= 1 \end{align*}
		Thay $x=1$ vào $y = 3x-1$:
		$y = 3(1) - 1 = 3 - 1 = 2$.
		Vậy, nghiệm của hệ phương trình là $(x;y) = (1;2)$.
		(Hoặc có thể giải bằng phương pháp cộng đại số)
	}
\end{ex}
\begin{ex}
	Một chiếc hộp có chứa 5 tấm thẻ cùng loại, được đánh số lần lượt là 3; 5; 6; 7; 9. Lấy ngẫu nhiên đồng thời 2 tấm thẻ từ hộp. Tính xác suất của mỗi biến cố: A: “Tổng các số ghi trên 2 tấm thẻ lớn hơn 13”.
	\loigiai{
		Các tấm thẻ được đánh số là $\{3, 5, 6, 7, 9\}$.
		Khi lấy ngẫu nhiên đồng thời 2 tấm thẻ từ hộp, ta sử dụng phương pháp liệt kê để tìm không gian mẫu (tất cả các cặp thẻ có thể lấy ra).
		Không gian mẫu $\Omega$ bao gồm các cặp (không phân biệt thứ tự):
		\begin{itemize}
			\item (3,5); (3,6); (3,7); (3,9)
			\item (5,6); (5,7); (5,9)
			\item (6,7); (6,9)
			\item (7,9)
		\end{itemize}
		Số phần tử của không gian mẫu là $n(\Omega) = 4+3+2+1 = 10$.
		
		Biến cố A: “Tổng các số ghi trên 2 tấm thẻ lớn hơn 13”.
	
		Các kết quả thuận lợi cho biến cố A là: $\{(5,9), (6,9), (7,9)\}$.
		Số phần tử của biến cố A là $n(A) = 3$.
		
		Xác suất của biến cố A là:
		$P(A) = \dfrac{n(A)}{n(\Omega)} = \dfrac{3}{10}$.
		
		Vậy, xác suất để tổng các số ghi trên 2 tấm thẻ lớn hơn 13 là $\dfrac{3}{10}$.
	}
\end{ex}
\begin{ex}
	Một phòng họp có 360 ghế ngồi được xếp thành từng dãy và số ghế của từng dãy đều như sau. Nếu giữ nguyên số ghế ngồi trong phòng nhưng giảm đi 2 dãy ghế thì số ghế trong mỗi dãy tăng thêm 6 ghế. Hỏi ban đầu trong phòng họp có bao nhiêu dãy ghế?
	\loigiai{
		Gọi số dãy ghế ban đầu trong phòng họp là $x$ (dãy).
		Điều kiện: $x$ là số nguyên dương.
		Khi đó, số ghế trong mỗi dãy ban đầu là $\dfrac{360}{x}$ (ghế).
		
		Nếu giảm đi 2 dãy ghế, số dãy ghế còn lại là $x-2$ (dãy).
		Điều kiện: $x-2 > 0 \Rightarrow x > 2$.
		Khi đó, số ghế trong mỗi dãy lúc sau là $\dfrac{360}{x-2}$ (ghế).
		
		Theo đề bài, nếu giảm đi 2 dãy ghế thì số ghế trong mỗi dãy tăng thêm 6 ghế, nên ta có phương trình:
		$\dfrac{360}{x-2} - \dfrac{360}{x} = 6$
		
		Để giải phương trình, ta quy đồng mẫu số:
		\begin{align*} \dfrac{360x - 360(x-2)}{x(x-2)} &= 6 \\ \dfrac{360x - 360x + 720}{x^2 - 2x} &= 6 \\ \dfrac{720}{x^2 - 2x} &= 6 \end{align*}
		Với điều kiện $x > 2$ (suy ra $x \ne 0$ và $x-2 \ne 0$), ta có:
		\begin{align*} 720 &= 6(x^2 - 2x) \\ 720 &= 6x^2 - 12x \\ 6x^2 - 12x - 720 &= 0 \end{align*}
		Chia cả hai vế cho 6:
		$x^2 - 2x - 120 = 0$
		
		Giải phương trình bậc hai này:
		$\Delta' = (-1)^2 - 1 \cdot (-120) = 1 + 120 = 121$
		$\sqrt{\Delta'} = \sqrt{121} = 11$
		
		Phương trình có hai nghiệm phân biệt:
		$x_1 = \dfrac{-(-1) + 11}{1} = 1 + 11 = 12$ (Thỏa mãn điều kiện $x > 2$)
		$x_2 = \dfrac{-(-1) - 11}{1} = 1 - 11 = -10$ (Không thỏa mãn điều kiện $x > 2$)
		
		Vậy, ban đầu trong phòng họp có 12 dãy ghế.
		
		\textbf{Thử lại:}
		Ban đầu: 12 dãy ghế, mỗi dãy có $360/12 = 30$ ghế.
		Lúc sau: $12-2 = 10$ dãy ghế, mỗi dãy có $360/10 = 36$ ghế.
		Số ghế mỗi dãy tăng thêm: $36 - 30 = 6$ ghế (Đúng).
	}
\end{ex}
\begin{ex}
	Cho tam giác ABC nhọn. Đường tròn (O) đường kính BC cắt AB, AC lần lượt tại E và D; BD cắt CE tại H, AH cắt BC tại I. Từ A kẻ tiếp tuyến AM, AN của đường tròn (O) (M, N là tiếp điểm).
	\begin{enumerate}
		\item[a)] Chứng minh 4 điểm A, E, H, D nằm trên một đường tròn.
		\item[b)] Chứng minh AB.BE = BI.BC từ đó suy ra AB.BE + AC.CD = BC$^2$.
	\end{enumerate}
\end{ex}

\begin{ex}
	\begin{enumerate}
		\item[a)] Người ta cắt bỏ bốn hình vuông cùng kích thước ở bốn góc của một tấm tôn hình vuông có cạnh 1 m để gò lại thành một chiếc thùng có dạng hình hộp chữ nhật không nắp. Hỏi cạnh của các hình vuông cần bỏ đi có độ dài bằng bao nhiêu để thùng hình hộp nhận được có thể tích lớn nhất?
		\item[b)] Cho a, b, c là các số thực dương thoả mãn: $a+b+c=3$.
		Tìm giá trị nhỏ nhất của biểu thức: $P = \dfrac{2025}{\sqrt{ab}+\sqrt{bc}+\sqrt{ac}} - \dfrac{b\sqrt{a}}{1+b} - \dfrac{c\sqrt{b}}{1+c} - \dfrac{a\sqrt{c}}{1+a}$.
	\end{enumerate}
\end{ex}

\setcounterref{ex}{0}
